\documentclass[11pt,a4paper]{article}
\usepackage[margin=2.2cm]{geometry}
\usepackage[T1]{fontenc}
\usepackage{lmodern}
\usepackage{amsmath,amssymb}
\usepackage{graphicx}
\usepackage{booktabs}
\usepackage{array}
\usepackage{longtable}
\usepackage{caption}
\usepackage{float}
\usepackage[dvipsnames]{xcolor}
\usepackage[most]{tcolorbox}
\usepackage{url}
\def\UrlBreaks{\do\/\do-\do_\do.}
\usepackage[colorlinks=true,linkcolor=blue!50!black,urlcolor=blue!50!black,citecolor=blue!50!black]{hyperref}
\DeclareGraphicsExtensions{.pdf,.png}
\graphicspath{{figures/}}
\sloppy
\emergencystretch=2em
\captionsetup{font=small,labelfont=bf}

% ---------------------------------------------------------------- macros
\newcommand{\fpath}[1]{\nolinkurl{#1}}
\newcommand{\fcmd}[1]{\texttt{\detokenize{#1}}}
\newcommand{\Ks}{K$^0_{\rm S}$}
\newcommand{\Lam}{$\Lambda$}
\newcommand{\aLam}{$\bar{\Lambda}$}
\newcommand{\Dzero}{D$^0$}
\newcommand{\Xc}{$X$}
\newcommand{\pt}{\ensuremath{p_T}}
\newtcolorbox{choicebox}{colback=blue!5,colframe=blue!60!black,fonttitle=\bfseries,title=Choice,left=2pt,right=2pt,top=1pt,bottom=1pt}
\newtcolorbox{caveatbox}{colback=red!5,colframe=red!60!black,fonttitle=\bfseries,title=Caveat,left=2pt,right=2pt,top=1pt,bottom=1pt}
\newtcolorbox{decisionbox}{colback=orange!7,colframe=orange!80!black,fonttitle=\bfseries,title=Decision needed,left=2pt,right=2pt,top=1pt,bottom=1pt}
\newcommand{\choice}[1]{\begin{choicebox}\small #1\end{choicebox}}
\newcommand{\caveat}[1]{\begin{caveatbox}\small #1\end{caveatbox}}
\newcommand{\decision}[1]{\begin{decisionbox}\small #1\end{decisionbox}}

% script-generated numbers (scripts/make_numbers.C, scripts/job_stats.py)
\input{tables/numbers.tex}
\input{tables/jobstats.tex}
\input{tables/trim_numbers.tex}

\title{KFParticle reconstruction of \Ks$\to\pi^+\pi^-$, \Lam$\to p\pi^-$ and \Dzero$\to K^-\pi^+$\\
in sPHENIX Run-3 p+p \texttt{DST\_TRKR\_TRACKS} (ana581\_2026p005\_v001)\\[4pt]
\large Analysis note: setup, pre-selection, ntuple content and first test}
\author{X.~Dong (dongxs) \quad -- \quad draft prepared with Hermes}
\date{\today}

\begin{document}
\maketitle

\begin{abstract}
\noindent
This note documents a KFParticle-based reconstruction of \Ks, \Lam(\aLam) and \Dzero(\={D}$^0$)
candidates on the Run-3 p+p track DSTs, with loose topological pre-cuts and all topological
variables written to flat ntuples. It lists every event, track, pair, vertex and mother-level
requirement, how the primary vertex (PV) and the bunch crossing of each track enter the
reconstruction, the ntuple content and its size, and the result of a first test on
\nEvents{} events of run 79509 (31 DST segments). With offline QA cuts the \Ks{} peak is at
$\meanKs\pm\emeanKs$\,MeV with $\sigma=\sigmaKs$\,MeV ($\yieldKs\pm\eyieldKs$ signal) and
the \Lam+\aLam{} peak at $\meanLambda\pm\emeanLambda$\,MeV with $\sigma=\sigmaLambda$\,MeV
($\yieldLambda\pm\eyieldLambda$). A dedicated study without any crossing requirement shows
that true \Ks{} and \Lam{} are found with both daughters and the PV in the same crossing,
supporting the default crossing-matched reconstruction; candidates with a TPC-only daughter are
only recoverable at crossing 0.
For mass production a trimmed configuration (daughter impact-parameter cuts, DIRA$>0.9$, slim
format without trigger, redundant momentum and covariance branches) reduces the candidate trees from
\trimKBLoose{} to \trimKBDira\,kB per event without a measurable change of the \Ks{} and \Lam{}
signal after the QA cuts.
\textbf{Limitations:} one run, one test sample, no event or trigger selection, no
efficiency correction, no systematic uncertainties; \Dzero{} shows no peak at this statistics.
Boxes: \textcolor{blue!60!black}{blue} = choice made and why,
\textcolor{red!60!black}{red} = caveat / open point,
\textcolor{orange!80!black}{orange} = needs a decision by the analyser.
\end{abstract}

\tableofcontents
\clearpage

\section{Overview}
\label{sec:overview}

\paragraph{Goal.} Prepare a reconstruction of the neutral strange hadrons \Ks{} and \Lam{} and
of the \Dzero{} meson from the Run-3 p+p track DSTs, keep the topological pre-selection loose,
store all topological information in a flat ntuple so that the final selection can be tuned
offline, and verify on a test sample that the reconstruction behaves as expected, in
particular with respect to the bunch-crossing assignment of each track in streaming readout.

\paragraph{Main findings of the test (\nEvents{} events, run 79509).}
\begin{itemize}\setlength\itemsep{1pt}
\item The ntuples contain \candPerEvKs{} \Ks, \candPerEvLambda{} \Lam{} and \candPerEvDzero{}
  \Dzero{} candidates per event; \nBranchKs{} branches; $\approx$\BPerCandKs{} bytes per
  candidate on disk. Together \kBPerEvProd{}\,kB per event (Section~\ref{sec:ntuple}).
\item Clear \Ks{} and \Lam{} peaks after offline QA cuts: \Ks{} $\mu=\meanKs\pm\emeanKs$\,MeV
  ($\massShiftKs$\,MeV from the PDG value), $\sigma=\sigmaKs$\,MeV, S/B=\SBKs;
  \Lam{} $\mu=\meanLambda\pm\emeanLambda$\,MeV, $\sigma=\sigmaLambda$\,MeV, S/B=\SBLambda{}
  (Section~\ref{sec:v0}).
\item Crossing logic: in the default sample every candidate has both daughters and the PV in the
  same crossing (\checkXKs/\checkXLambda). Without any crossing requirement the \Ks{} signal
  appears only in the same-crossing classes; for \Lam{} the yield in the
  ``daughters in different crossings'' class is $\xclCLambda\pm\exclCLambda$, i.e.\ not
  significant (Section~\ref{sec:crossing}).
\item TPC-only tracks carry crossing 0 by construction (\trkTPConlyFrac\% of all tracks); they
  can only form candidates with the vertex of crossing 0, which is present in
  \evVtxXzeroFrac\% of the events.
\item The same collision can be reconstructed in two consecutive streaming events:
  \dupFracKs\% (\Ks) and \dupFracLambda\% (\Lam) of the candidates are such duplicates
  (identified by \fcmd{Collision_BCO}).
\item No \Dzero{} peak is visible in the test sample (Section~\ref{sec:d0}).
\item Production configuration (Section~\ref{sec:trim}): \trimKBDira\,kB per event for the three
  channels instead of \trimKBLoose{} (plus \trimEvInfoKB\,kB for \fcmd{EventInfo}), $\approx$\trimTBDira\,TB
  for the full dataset; \Ks{} and \Lam{} QA-cut signal unchanged.
\end{itemize}

\paragraph{Choices in one line each.}
\begin{itemize}\setlength\itemsep{1pt}
\item Release ana.581, identical to the DST production (Section~\ref{sec:dataset}).
\item No event-level selection (no trigger, vertex-$z$ or run-quality cut) at ntuple level.
\item Primary vertices from \fcmd{SvtxVertexMap} (track vertices), not GlobalVertex/MBD.
\item Daughters and PV must share one bunch crossing (default KFParticle behaviour, kept).
\item Loose pre-cuts: no lifetime, decay-length, flight-distance or impact-parameter
  requirements; DIRA$>0.8$ (V0) / $>0$ (\Dzero), DCA$_{12}<5$\,mm (V0) / 500\,$\mu$m (\Dzero),
  vertex $\chi^2/{\rm ndf}<20$ (Table~\ref{tab:precuts}).
\item No PID requirement; TPC d$E$/d$x$ stored for offline use.
\item \Dzero{} daughters need $\geq2$ MVTX clusters and $p_T>0.2$\,GeV/$c$.
\item Per-cluster detector vectors and all-PV vectors switched off to limit the ntuple size.
\item A per-event companion tree records all processed events and the full vertex list.
\item Production cut set: daughter DCA to PV $>0.05$\,cm (V0) / $>30\,\mu$m (\Dzero), DIRA$>0.9$, no
  crossing-study output, slim format without trigger, redundant momentum and covariance branches.
\end{itemize}

\section{Dataset, software and computing}
\label{sec:dataset}

\paragraph{Input.} Official Run-3 p+p track DSTs,
\fpath{/sphenix/lustre01/sphnxpro/production/run3pp/physics/ana581_2026p005_v001/DST_TRKR_TRACKS/}.
At the time of the test the area contained \nDSTfiles{} files in four run ranges
(\fpath{logs/dataset_inventory.txt}); the production was still growing. Each segment holds
$\approx$\evPerFile{} events. The DST nodes and their compressed size per event (segment
79509-00000, \fpath{logs/inspect_dst_79509_00000.txt}) are: \fcmd{TRKR_CLUSTER_SEED} 144\,kB,
\fcmd{SiliconTrackSeedContainer} 41\,kB, \fcmd{SvtxTrackMap} (\fcmd{SvtxTrack_v5}) 13\,kB,
\fcmd{TpcTrackSeedContainer} 11\,kB, \fcmd{SvtxTrackSeedContainer} 4\,kB,
\fcmd{SvtxVertexMap} 0.6\,kB, \fcmd{GL1RAWHIT}, \fcmd{EventHeader}, \fcmd{Sync}. There is no
MBD vertex, no calorimeter information and no truth information.

\paragraph{Test sample.} Run 79509, segments 00000--00030 (31 files,
\fpath{logs/dst_list.txt}), \nEvents{} events. Segment 0 was processed interactively on
sphnxuser01, segments 1--30 as HTCondor cluster 471173 (\nCondorOK/\nCondorJobs{} jobs with exit
code 0, \fpath{logs/condor_471173.txt}); all \nKFPdone/\nKFPlogs{} job logs end with the macro's
``All done'' line. Merged ntuples:
\fpath{/gpfs02/sphenix/user/dongxs/run3pp/output/merged/KFP_<channel>_run79509_seg000-030.root},
local copies in \fpath{~/Work/Hermes/sPHENIX/run3pp/ntuples/} (md5 verified).

\caveat{A single run is used. Its trigger mix, luminosity and run-quality status were not
checked. Two further single-segment tests (runs 80906 and 81600) were still running when this
note was written and are not included.}

\paragraph{Software.} sPHENIX release ana.581 (\fpath{logs/software.txt}):
coresoftware \fcmd{3ae2baa290a3}, macros \fcmd{87f5cb5e2ba1}, acts \fcmd{2f27de2d0f02},
calibrations \fcmd{476e33a37dc2}; ROOT 6.32.06, gcc 14.2.0. The KFParticle\_sPHENIX package is the
one of this coresoftware commit; statements in this note about its internal logic were read from
that source.

\choice{ana.581 is used because the DSTs were produced with ana581; this keeps the track and
vertex classes, the CDB tag logic and the field map consistent with the production. The
environment is set by \fpath{setup_env.sh} (release + local install only; the login environment
with an older private install is not used).}

\paragraph{Code} (cluster: \fpath{/gpfs02/sphenix/user/dongxs/run3pp}, local mirror
\fpath{~/Work/Hermes/sPHENIX/run3pp}):
\begin{itemize}\setlength\itemsep{0pt}
\item \fpath{macro/Fun4All_run3pp_KFP.C} -- Fun4All macro (all cuts of Section~\ref{sec:reco});
  \fpath{macro/run_kfp.sh} -- wrapper (log, ``All done'' check).
\item \fpath{src/EventInfoTree/} -- per-event companion module (library \fcmd{libeventinfotree}).
\item \fpath{condor/} -- file lists and submission (\fcmd{submit.sh} refuses duplicate submissions).
\item \fpath{analysis_note/scripts/make_numbers.C} -- every number, table and figure of this
  note; \fpath{analysis_note/scripts/job_stats.py} -- job statistics.
\end{itemize}

\paragraph{Computing.} Per-event time of the KFParticle modules (mean over the
\nKFPlogs{} jobs): \Ks{} \msPerEvKs\,ms, \Lam{} \msPerEvLambda\,ms, \Dzero{} \msPerEvDzero\,ms,
i.e.\ \msPerEvProd\,ms for the three default channels ($\approx$\minPerFileProd\,min per DST
segment); the two crossing-study channels add \msPerEvKsNoX{}+\msPerEvLambdaNoX\,ms. Condor
wall time with all five channels: \condorWallMin\,min on average (max \condorWallMaxMin\,min);
maximum resident memory \maxRSSGB\,GB.

\section{Reconstruction and pre-selection}
\label{sec:reco}

\subsection{Fun4All chain}
For every DST segment the macro \fpath{macro/Fun4All_run3pp_KFP.C} sets the run number and
timestamp from the file name, uses the CDB global tag \fcmd{newcdbtag}, loads the tracking
geometry (\fcmd{Tracking_Geometry} from the CDB) through a run-node input manager, calls
\fcmd{TpcReadoutInit(run)}, enables MVTX misalignment and \fcmd{TRACKING::streaming_mode},
and runs \fcmd{TrackingInit()} (ACTS geometry, field map \fcmd{FIELDMAP_TRACKING}). This mirrors the
official example \fpath{macros/TrackingProduction/Fun4All_KShortReco.C} of the same macros commit.
No track or vertex is refitted: the analysis uses the production \fcmd{SvtxTrackMap} and
\fcmd{SvtxVertexMap} as they are. The registered modules are, in order:
\fcmd{TriggerRunInfoReco} (trigger names for the trigger-bit branches), \fcmd{EventInfoTree},
and one \fcmd{KFParticle_sPHENIX} instance per channel (\fcmd{Ks_reco}, \fcmd{Lambda_reco},
\fcmd{D0_reco}; with \fcmd{crossingStudy=true} also \fcmd{Ks_noXmatch}, \fcmd{Lambda_noXmatch}).
Each instance writes its own file \fcmd{KFP_<channel>_<run>_<segment>_<skip>.root}.

\subsection{Event-level selection}
\choice{None. Every event of the DST is processed and recorded in the \fcmd{EventInfo} tree,
candidates are written for all events. No trigger, vertex-$z$, pile-up or run-quality cut is
applied at ntuple level; trigger bits (64 \fcmd{trigger_*} branches per candidate, and
\fcmd{GL1_scaled}/\fcmd{GL1_live} per event) and all vertices are stored so these cuts can be
applied offline.}

In the test sample an event (one GL1 trigger in streaming readout) contains on average
\trkPerEv{} tracks and \vtxPerEv{} reconstructed vertices spread over beam crossings
\vtxXmin{} to \vtxXmax{} relative to the GL1 BCO (Figure~\ref{fig:event}).

\begin{figure}[H]
\centering
\includegraphics[width=\textwidth]{fig_event_overview}
\caption{Event overview of the test sample (\nEvents{} events, \fcmd{EventInfo} tree, no
selection). Left: beam crossing of all \fcmd{SvtxVertexMap} vertices relative to the event's GL1
BCO (log scale); the spike at 0 holds \vtxAtXzero{} vertices compared with
$\approx$\vtxPerOtherX{} per crossing elsewhere. Middle: vertex $z$. Right: number of tracks per
event (black) and of TPC-only tracks per event (red).}
\label{fig:event}
\end{figure}

\subsection{Track classes and bunch crossing}
\label{sec:trackclasses}
Each track in \fcmd{SvtxTrackMap} carries a bunch-crossing number \Xc{} relative to the GL1 BCO
(\fcmd{SvtxTrack::get_crossing()}); each vertex carries the crossing of its tracks
(\fcmd{get_beam_crossing()}). We classify tracks by their silicon clusters
(\fcmd{SvtxTrack_v5} cluster counts):
\begin{itemize}\setlength\itemsep{0pt}
\item \textbf{I}: $\geq1$ INTT cluster -- crossing from INTT timing (\trkINTTPerEv{} per event);
\item \textbf{M}: MVTX but no INTT clusters (\trkMVTXonlyPerEv{} per event); these tracks have
  non-zero crossings in the DST, but how their crossing is determined was not checked;
\item \textbf{T}: TPC-only, no silicon cluster (\trkTPConlyPerEv{} per event, \trkTPConlyFrac\%).
  All TPC-only tracks have $X=0$ in the DST (17871 of 17871 in segment 0,
  \fpath{logs/inspect_dst_79509_00000.txt}): their crossing is not measured, and their
  $z$ position is only correct if they really come from crossing 0.
\end{itemize}

\subsection{Track (daughter) pre-selection}
\label{sec:trackcuts}
All cuts are applied by KFParticle\_sPHENIX. Cluster counts (``hits'' in the KFParticle
interface) are the \fcmd{SvtxTrack_v5} cluster numbers. The track quality cuts (\pt, $\chi^2$, DCA
to PV) are evaluated for each track pair after the pair DCA and two-track vertex $\chi^2$ cuts.

\begin{table}[H]
\centering\small
\caption{Daughter-track requirements. ``--'' = no requirement (set to $\pm$FLT\_MAX or the
KFParticle default that does not cut).}
\label{tab:trackcuts}
\begin{tabular}{lccc}
\toprule
 & \Ks$\to\pi^+\pi^-$ & \Lam$\to p\pi^-$ (+c.c.) & \Dzero$\to K^-\pi^+$ (+c.c.) \\
\midrule
TPC clusters & $\geq20$ & $\geq20$ & $\geq20$ \\
MVTX clusters & $\geq0$ & $\geq0$ & $\geq2$ \\
INTT clusters & $\geq0$ & $\geq0$ & $\geq0$ \\
TPOT clusters & $\geq0$ & $\geq0$ & $\geq0$ \\
\pt{} (GeV/$c$) & $\geq0.1$ & $\geq0.1$ & $\geq0.2$ \\
track $\chi^2$/ndf & -- & -- & -- \\
DCA to PV (3D, and its significance) & -- & -- & -- \\
DCA$_{xy}$ to PV & -- & -- & -- \\
PID (d$E$/d$x$) & none (stored) & none (stored) & none (stored) \\
mass hypothesis & $\pi$,$\pi$ & $p$,$\pi$ & $K$,$\pi$ \\
\bottomrule
\end{tabular}
\end{table}

\caveat{Implicit track requirement (read from \fcmd{KFParticle_Tools::isGoodTrack/calcMinPV_DCA},
not tested separately): the minimum DCA to the PVs is computed only with PVs of the track's own
crossing (for TPC-only tracks at $X=0$: with all PVs). A silicon track whose crossing has no
reconstructed vertex gets NaN, fails the ``DCA$\geq-1$'' comparison and is rejected. So every
daughter must have a vertex in its crossing -- also in the \fcmd{_noXmatch} instances.}

\subsection{Pair, vertex and mother pre-selection}
For each pair of selected tracks KFParticle (i) requires both tracks to have the same crossing
(\fcmd{requireBunchCrossingMatch}, default instances), (ii) cuts on the 3D distance of closest
approach of the two trajectories (DCA$_{12}$) and on its $xy$ projection, (iii) fits a two-track
vertex (KFVertex) and cuts on its $\chi^2$/ndf, (iv) applies the track cuts of
Table~\ref{tab:trackcuts}, (v) builds the mother for each allowed mass assignment and cuts on
mass, \pt{} and the volume of the mother's position-error ellipsoid, and (vi) for every PV
computes the pointing quantities w.r.t.\ that PV and requires the PV crossing to equal the
daughter crossing (\fcmd{requireTrackVertexBunchCrossingMatch}).

\begin{table}[H]
\centering\small
\caption{Pair, secondary-vertex, mother and PV-related pre-cuts as configured in
\fpath{macro/Fun4All_run3pp_KFP.C}. Lengths in cm.}
\label{tab:precuts}
\begin{tabular}{lccc}
\toprule
 & \Ks & \Lam & \Dzero \\
\midrule
daughters in same crossing & yes & yes & yes \\
DCA$_{12}$ (3D) & $<0.5$ & $<0.5$ & $<0.05$ \\
DCA$_{12}$ ($xy$) & -- & -- & -- \\
SV fit $\chi^2$/ndf & $<20$ & $<20$ & $<20$ \\
SV radius & -- & -- & -- \\
mother mass window (GeV/$c^2$) & 0.40--0.60 & 1.05--1.20 & 1.60--2.10 \\
mother \pt & -- & -- & -- \\
mother position-error ellipsoid volume (cm$^3$) & $<0.1$ & $<0.1$ & $<0.1$ \\
PV crossing = daughter crossing & yes & yes & yes \\
DIRA (3D) w.r.t.\ PV & $>0.8$ & $>0.8$ & $>0$ \\
DIRA$_{xy}$ & -- & -- & -- \\
flight-distance $\chi^2$ & -- & -- & -- \\
decay length, decay length$_{xy}$, $c\tau$ (and significances) & -- & -- & -- \\
mother DCA to PV (3D, $xy$, significances) & -- & -- & -- \\
charge conjugate & -- & included & included \\
\bottomrule
\end{tabular}
\end{table}

\choice{Lifetime and impact-parameter variables are not cut at ntuple level (they are stored),
so that the final selection, including a cut-based or MVA optimisation, can be done offline.
The kept cuts (DCA$_{12}$, SV $\chi^2$, DIRA, mass window) bound the combinatorics. For \Dzero{}
the much shorter $c\tau$ requires MVTX precision on both daughters, hence $\geq2$ MVTX clusters
and DCA$_{12}<500\,\mu$m.}

\subsection{Primary-vertex choice and candidate multiplicity}
The PVs are all vertices of \fcmd{SvtxVertexMap} (\fcmd{dontUseGlobalVertex(true)}; the MBD is not
used). For a given track pair every (PV, mass assignment) combination that passes all cuts is a
possible candidate. KFParticle keeps \emph{one} of them: the one whose mother has the smallest 3D
$\chi^2$-deviation from its PV (\fcmd{selectBestCombination}). Consequences:
\begin{itemize}\setlength\itemsep{0pt}
\item each track pair yields at most one candidate per channel;
\item for \Lam{} and \Dzero{} with the \fcmd{[...]cc} descriptor only one of the two mass
  assignments of a pair is kept. For \Dzero{} the swapped ($K\leftrightarrow\pi$) mass is
  recomputed offline from the stored daughter momenta (Section~\ref{sec:d0}); for \Lam{} the
  $p\pi^-$ vs $\bar p\pi^+$ ambiguity is resolved by charge, so only the $\pi\pi$-like
  ``wrong'' assignment would be lost, which is outside the mass window anyway.
\end{itemize}

\subsection{What is stored for the mother}
The mother mass and momentum are those of the unconstrained secondary-vertex fit (no mass or PV
constraint). Decay length, decay time and their $xy$ versions are computed from a copy of the
mother with its production vertex set to the chosen PV. DIRA, flight-distance $\chi^2$ and the
mother's DCA to the PV are computed w.r.t.\ the chosen PV. Daughter momenta are stored after the
daughters are moved to the decay vertex (\fcmd{extrapolateTracksToSV}, default true).

\subsection{Offline QA cuts used in this note}
\label{sec:qacuts}
To display and fit the V0 peaks we apply, offline only, DIRA$>0.995$, $L_{xy}>0.5$\,cm,
DCA$_{12}<0.2$\,cm and $|{\rm DCA}_{xy}|>0.02$\,cm to the PV for both daughters. These cuts were
chosen by eye to make the peaks visible; they are not optimised and are not part of the
ntuple pre-selection.

\section{Ntuple content and size}
\label{sec:ntuple}

\subsection{Candidate tree \fcmd{DecayTree}}
One file per channel and DST segment, one entry per candidate, \nBranchKs{} branches. Mother
prefix \fcmd{K_S0_}, \fcmd{Lambda0_}, \fcmd{D0_}; daughters \fcmd{track_1_}, \fcmd{track_2_}
(the PDG code tells which is which; for \Lam{} \fcmd{track_1} can be the pion or the proton).
\begin{table}[H]
\centering\small
\caption{Branch groups of \fcmd{DecayTree} (\fcmd{M} = mother prefix, \fcmd{d} = daughter
prefix). All lengths in cm, momenta in GeV/$c$.}
\label{tab:branches}
\begin{tabular}{p{0.24\textwidth}p{0.70\textwidth}}
\toprule
group & branches \\
\midrule
mother kinematics & \fcmd{M_mass, M_massErr, M_px/py/pz/pE/p/pErr, M_pT, M_pTErr, M_charge, M_pseudorapidity, M_rapidity, M_theta, M_phi, M_PDG_ID} \\
mother topology & \fcmd{M_x/y/z} (SV), \fcmd{M_decayLength(Err), M_decayLength_xy, M_decayLengthErr_xy, M_decayTime(Err), M_DIRA, M_DIRA_xy, M_FDchi2, M_PV_DCA, M_PV_DCA_xy, M_PV_DCA_Err, M_PV_DCA_StdDev, M_vertex_volume, M_chi2, M_nDoF, M_Covariance[21]} \\
pair & \fcmd{track_1_track_2_DCA, _DCA_xy, _DCA_sig, _DCA_sig_xy} \\
daughter kinematics & \fcmd{d_mass, d_px/py/pz/pE/p/pErr, d_pT, d_pTErr, d_jT, d_charge, d_pseudorapidity, d_rapidity, d_theta, d_phi, d_x/y/z, d_PDG_ID, d_track_ID, d_Covariance[21]} \\
daughter quality / topology & \fcmd{d_chi2, d_nDoF, d_PV_DCA, d_PV_DCA_Err, d_PV_DCA_xy, d_PV_DCA_sig, d_PV_DCA_sig_xy, d_bunch_crossing} \\
daughter detector & \fcmd{d_MVTX/INTT/TPC/TPOT_nHits} (clusters on the seeds), \fcmd{d_..._nStates}, \fcmd{d_calculated_dEdx} \\
PV & \fcmd{primary_vertex_x/y/z, _volume, _chi2, _nDoF, _Covariance[6], primary_vertex_ID} \\
event & \fcmd{runNumber, eventNumber, GL1_BCO, last_GL1_BCO, Collision_BCO, nPrimaryVerticesOfBC, nTracksOfBC, nTracksOfVertex, nSiSeedMultiplicity, secondary_vertex_mass_pionPID} \\
trigger & 64 Booleans \fcmd{trigger_<name>} (e.g.\ \fcmd{trigger_MBD_NandS_geq_1}) \\
\bottomrule
\end{tabular}
\end{table}

\caveat{Two branches do not mean what their names suggest:
\fcmd{primary_vertex_ID} holds the \emph{beam crossing} of the chosen PV (KFParticle
\fcmd{getPVID} returns \fcmd{get_beam_crossing()}), not the vertex map key;
\fcmd{Collision_BCO} = \fcmd{GL1_BCO} + crossing of \fcmd{track_1}, i.e.\ the absolute BCO of the
collision, which is the right key to identify the same collision across events.}

\subsection{Per-event tree \fcmd{EventInfo}}
File \fcmd{KFP_EventInfo_*.root}, one entry per processed event (\nEvents{} entries in the test):
\fcmd{runNumber, eventNumber, GL1_BCO, GL1_scaled, GL1_live, GL1_bunch}; track counts \fcmd{nTrk,
nTrkINTT, nTrkMVTXonly, nTrkTPConly, nTrkX0}; and the vertex list \fcmd{nVtx, vtx_id,
vtx_crossing, vtx_x/y/z, vtx_chi2, vtx_ndof, vtx_ntrk}. It is used for normalisation and to look
up any vertex by (\fcmd{runNumber}, \fcmd{eventNumber}, crossing).

\subsection{Size}
\begin{table}[H]
\centering\small
\caption{Ntuple size in the test sample (\nEvents{} events; \fpath{tables/tab_sizes.tex} from
\fcmd{TTree::GetZipBytes/GetTotBytes}).}
\label{tab:sizes}
\begin{tabular}{lrrrrrrr}
\toprule
channel & candidates & per event & branches & B/cand (raw) & B/cand (disk) & kB/event & MB total \\
\midrule
\input{tables/tab_sizes.tex}
\bottomrule
\end{tabular}
\end{table}
Together the three channels need \kBPerEvProd\,kB per event. Extrapolated to the \nDSTfiles{}
segments available at the time of the test ($\approx$\nEvFullE$\times10^8$ events) this is
$\approx$\TBFull\,TB and $\approx$\kCPUhFull{}k CPU hours for the three default channels
(\fpath{logs/job_stats.txt}).

\choice{These sizes are for the loose diagnostic configuration. For mass production a trimmed
configuration is used (Section~\ref{sec:trim}): \trimKBDira\,kB per event for the three channels.}

\section{\Ks{} and \Lam{} in the default (crossing-matched) sample}
\label{sec:v0}

Figures~\ref{fig:precutKs} and \ref{fig:precutLambda} show the stored topological variables,
i.e.\ after the reconstruction pre-cuts and before any offline cut. Figures~\ref{fig:massKs}
and~\ref{fig:massLambda} show the invariant mass before and after the offline QA cuts of
Section~\ref{sec:qacuts}, after removing duplicates (Section~\ref{sec:dup}). Fits: Gaussian plus
second-order polynomial, binned likelihood, in 0.43--0.57 (\Ks) and 1.09--1.16\,GeV/$c^2$ (\Lam);
the signal yield is the Gaussian integral, S/B is in $\pm3\sigma$. Uncertainties are statistical.

\begin{figure}[H]
\centering
\includegraphics[width=\textwidth]{fig_precut_Ks}
\caption{\Ks{} candidates as stored (pre-cuts only, \nCandKs{} candidates, \nEvents{} events):
DIRA, DCA between the daughters, SV $\chi^2$/ndf, decay length, the smaller of the two daughter
DCA$_{xy}$ to the PV, and mother \pt. Red dashed lines: pre-cut values (DIRA$>0.8$,
DCA$_{12}<0.5$\,cm, $\chi^2$/ndf$<20$). Log scale.}
\label{fig:precutKs}
\end{figure}

\begin{figure}[H]
\centering
\includegraphics[width=\textwidth]{fig_precut_Lambda}
\caption{As Figure~\ref{fig:precutKs} for \Lam+\aLam{} candidates (\nCandLambda{} candidates).}
\label{fig:precutLambda}
\end{figure}

\begin{figure}[H]
\centering
\includegraphics[width=0.9\textwidth]{fig_mass_Ks}
\caption{$\pi^+\pi^-$ invariant mass. Left: pre-cuts only. Right: with offline QA cuts
(\nQAKs{} candidates before, \dupQAKs{} duplicates removed) and the fit (red). The \Ks{} peak is
at $\meanKs\pm\emeanKs$\,MeV, $\massShiftKs$\,MeV below the PDG mass (497.6\,MeV), with
$\sigma=\sigmaKs\pm\esigmaKs$\,MeV, $N_{\rm sig}=\yieldKs\pm\eyieldKs$ and S/B=\SBKs.}
\label{fig:massKs}
\end{figure}

\begin{figure}[H]
\centering
\includegraphics[width=0.9\textwidth]{fig_mass_Lambda}
\caption{$p\pi$ invariant mass (\Lam{} and \aLam{} together). Left: pre-cuts only. Right: with
offline QA cuts and fit: $\mu=\meanLambda\pm\emeanLambda$\,MeV, $\sigma=\sigmaLambda\pm\esigmaLambda$\,MeV,
$N_{\rm sig}=\yieldLambda\pm\eyieldLambda$, S/B=\SBLambda.}
\label{fig:massLambda}
\end{figure}

With the shape fixed to the inclusive fit, the separate yields are
\Lam: $\yieldLambdaP\pm\eyieldLambdaP$ and \aLam: $\yieldLambdaA\pm\eyieldLambdaA$.
The QA-cut signal per 1000 events is \yieldPerKevKs{} (\Ks) and \yieldPerKevLambda{} (\Lam+\aLam);
these are not efficiency corrected and depend strongly on the ad-hoc QA cuts.

\caveat{The \Ks{} mass is $\massShiftKs$\,MeV below the PDG value, i.e.\ about
$-0.6\%$ in mass; with the stated statistical error this is significant. Possible causes
(momentum scale / field map, energy-loss correction, daughter momenta taken at the SV) were not
investigated. The \Lam{} mass agrees with the PDG value within $\approx$1\,MeV.}

\subsection{Yield by daughter type}
\begin{table}[H]
\centering\small
\caption{Signal by daughter-track type (I = INTT, M = MVTX-only, T = TPC-only;
Section~\ref{sec:trackclasses}), default sample with QA cuts, duplicates removed. Gaussian mean
and width fixed to the inclusive fit, polynomial background free.}
\label{tab:pairtype}
\begin{minipage}[t]{0.48\textwidth}\centering
\begin{tabular}{lrrr}
\toprule
\multicolumn{4}{c}{\Ks} \\
types & cand. & signal & \% \\
\midrule
\input{tables/tab_pairtype_Ks.tex}
\bottomrule
\end{tabular}
\end{minipage}\hfill
\begin{minipage}[t]{0.48\textwidth}\centering
\begin{tabular}{lrrr}
\toprule
\multicolumn{4}{c}{\Lam+\aLam} \\
types & cand. & signal & \% \\
\midrule
\input{tables/tab_pairtype_Lambda.tex}
\bottomrule
\end{tabular}
\end{minipage}
\end{table}
Most of the signal has two INTT-matched daughters. Pairs with two TPC-only daughters (TT) carry
a small signal; by construction they come with a crossing-0 PV (Section~\ref{sec:crossing}).
Mass distributions per type: Figures~\ref{fig:pairKs}--\ref{fig:pairLambda}.

\begin{figure}[H]
\centering
\includegraphics[width=0.95\textwidth]{fig_pairtype_Ks}
\caption{$\pi^+\pi^-$ mass by daughter type (default sample, QA cuts, duplicates removed); blue:
fit with mean and width fixed to the inclusive fit. Numbers in Table~\ref{tab:pairtype}.}
\label{fig:pairKs}
\end{figure}
\begin{figure}[H]
\centering
\includegraphics[width=0.95\textwidth]{fig_pairtype_Lambda}
\caption{As Figure~\ref{fig:pairKs} for $p\pi$.}
\label{fig:pairLambda}
\end{figure}

\section{Bunch-crossing assignment: checks}
\label{sec:crossing}

\subsection{Why the crossing matters}
In streaming readout one event (one GL1 trigger) contains tracks from many collisions in different
bunch crossings (\vtxPerEv{} vertices per event, Figure~\ref{fig:event}). The silicon detectors
fix the crossing of a track (INTT timing); the TPC measures only the drift time, so the $z$
position of the TPC part of a track depends on the crossing assumed for it. Combining two tracks
from different crossings, or attaching a candidate to a PV from another crossing, produces
combinatorial background and wrong pointing/decay-length values. The default reconstruction
therefore requires both daughters and the PV to have the same crossing
(Section~\ref{sec:reco}).

The vertex-crossing distribution (Figure~\ref{fig:event}, left) is flat at
$\approx$\vtxPerOtherX{} vertices per crossing, with a spike of \vtxAtXzero{} vertices at $X=0$;
\evVtxXzeroFrac\% of the events have a vertex at $X=0$ (\evVtxXzeroMultiFrac\% more than one).
We interpret the spike as the collision that fired the trigger; this interpretation was not
checked against the trigger information. Because all TPC-only tracks have $X=0$, a candidate
with a TPC-only daughter can only be formed with $X=0$ silicon tracks and an $X=0$ PV.

\subsection{Consistency of the default sample}
In the default \Ks{} and \Lam{} ntuples every candidate satisfies
\fcmd{track_1_bunch_crossing} = \fcmd{track_2_bunch_crossing} = \fcmd{primary_vertex_ID}
(violations: \nViolXKs{} and \nViolXLambda; check \checkXKs/\checkXLambda,
\fpath{scripts/make_numbers.C}). \nTTKs{} \Ks{} and \nTTLambda{} \Lam{} candidates have two
TPC-only daughters (before QA cuts).

\subsection{Duplicates across events}
\label{sec:dup}
Consecutive streaming events overlap in time, so the same collision can be reconstructed twice,
with crossings that differ by the BCO difference of the two events. Such candidates have the same
\fcmd{Collision_BCO} and identical kinematics. Using the key (\fcmd{Collision_BCO}, mass rounded to
10\,keV, $p_x$ of \fcmd{track_1} rounded to 0.1\,MeV) we find \dupFracKs\% (\Ks) and
\dupFracLambda\% (\Lam) duplicates before QA cuts, \dupQAKs{} and \dupQALambda{} after QA cuts. They
are removed in all results of this note.

\choice{Duplicates are kept in the ntuple and removed offline with the \fcmd{Collision_BCO} key: the
removal needs the neighbouring events, which a per-event module cannot see across file
boundaries.}

\subsection{Reconstruction without any crossing requirement}
To test whether the crossing requirement loses real signal, the \fcmd{_noXmatch} instances run with
the same cuts but with \fcmd{requireBunchCrossingMatch(false)} and
\fcmd{requireTrackVertexBunchCrossingMatch(false)}: any two tracks may be combined and the PV is
chosen among all vertices by the smallest $\chi^2$-deviation. (The implicit requirement that each
daughter has a vertex in its own crossing, Section~\ref{sec:trackcuts}, still applies.) This gives
\candPerEvNoXKs{} \Ks{} and \candPerEvNoXLambda{} \Lam{} candidates per event. Each candidate is
put in one of six classes:
\begin{itemize}\setlength\itemsep{0pt}
\item[A] both daughters with silicon, $X_1=X_2=X_{\rm PV}$ (what the default keeps);
\item[B] both with silicon, $X_1=X_2\neq X_{\rm PV}$ (consistent pair, PV of another crossing);
\item[C] both with silicon, $X_1\neq X_2$;
\item[D] one TPC-only daughter ($X=0$) and one silicon daughter with $X\neq0$;
\item[E] one TPC-only daughter and one silicon daughter with $X=0$;
\item[F] two TPC-only daughters.
\end{itemize}
The signal in each class is obtained with the QA cuts and a fit with mean and width fixed to the
default-sample result (Table~\ref{tab:xclass}, Figures~\ref{fig:xclassKs}--\ref{fig:xclassLambda}).

\begin{table}[H]
\centering\small
\caption{Signal per crossing class in the no-crossing-requirement samples (QA cuts; Gaussian
with fixed mean/width + free pol2; statistical errors). Duplicates are not removed here.}
\label{tab:xclass}
\begin{tabular}{llrrr}
\toprule
\multicolumn{5}{c}{\Ks} \\
class & definition & candidates & signal & signal/error \\
\midrule
\input{tables/tab_xclass_Ks.tex}
\midrule
\multicolumn{5}{c}{\Lam+\aLam} \\
class & definition & candidates & signal & signal/error \\
\midrule
\input{tables/tab_xclass_Lambda.tex}
\bottomrule
\end{tabular}
\end{table}

\paragraph{Observations.}
\begin{itemize}\setlength\itemsep{1pt}
\item \textbf{Class C} (daughters in different crossings) contains no significant signal:
  \Ks{} $\xclCKs\pm\exclCKs$, \Lam{} $\xclCLambda\pm\exclCLambda$ ($2.6\sigma$), on a background
  of $>3\times10^4$ candidates. Splitting class C, the \Lam{} excess sits entirely at
  $|X_1-X_2|>5$ ($\xclClargeLambda\pm\exclClargeLambda$; $|X_1-X_2|\leq5$:
  $\xclCsmallLambda\pm\exclCsmallLambda$), where the background is largest; for \Ks{} both parts
  are compatible with zero ($\xclCsmallKs\pm\exclCsmallKs$, $\xclClargeKs\pm\exclClargeKs$).
\item A fit-free cross-check: for silicon--silicon pairs, peak window ($|m-\mu|<2\sigma$) minus
  equal-width sidebands (5--7$\sigma$) gives \Ks{} $\dXzeroAllKs\pm\edXzeroAllKs$ with
  $X_1=X_2$ and $\dXnonzeroAllKs\pm\edXnonzeroAllKs$ with $X_1\neq X_2$; for \Lam{}
  $\dXzeroAllLambda\pm\edXzeroAllLambda$ and $\dXnonzeroAllLambda\pm\edXnonzeroAllLambda$
  (Figure~\ref{fig:dX}, which shows $|X_1-X_2|\leq20$).
\item \textbf{Class B} holds $\xclBKs\pm\exclBKs$ \Ks: same-crossing pairs for which the
  smallest-$\chi^2$ PV is from another crossing. Without the crossing requirement the PV choice
  is therefore wrong for about 10\% of the real \Ks{} (B/(A+B)), and the pointing variables are computed
  w.r.t.\ the wrong vertex. In the default sample these \Ks{} are attached to the PV of their own
  crossing: the default yields more \Ks{} with two silicon daughters (sum of the II, IM and MM rows of
  Table~\ref{tab:pairtype}) than class A ($\xclAKs\pm\exclAKs$), presumably because pointing cuts evaluated
  w.r.t.\ a wrong PV reject part of the signal.
\item \textbf{Class D} (TPC-only daughter + silicon daughter with $X\neq0$) has no signal
  (\Ks{} $\xclDKs\pm\exclDKs$, \Lam{} $\xclDLambda\pm\exclDLambda$): with the wrong crossing the
  TPC-only track does not form a good vertex. Such decays are lost both with and without the
  crossing requirement.
\item Classes E and F (crossing 0) carry small signals, consistent with the TT/IT entries of
  Table~\ref{tab:pairtype}.
\end{itemize}

\choice{Keep the default: both daughters and the PV in the same crossing. Relaxing it adds class C,
which after QA cuts has about 16 times more \Ks{} candidates than class A (Table~\ref{tab:xclass})
but no significant signal, and it makes the PV choice wrong for about 10\% of the \Ks{} (class B).}

\caveat{The \Lam{} class-C excess of $\xclCLambda\pm\exclCLambda$ is not significant but not zero
either. If real, it would be \Lam{} whose daughters were given different crossings (e.g.\ a
mis-assigned crossing of a low-momentum proton). This should be rechecked with $\sim$10 times more
events and with the crossing difference resolved per daughter type before it is ignored.}

\begin{figure}[H]
\centering
\includegraphics[width=0.48\textwidth]{fig_dX_Ks}\hfill
\includegraphics[width=0.48\textwidth]{fig_dX_Lambda}
\caption{Crossing difference $X_1-X_2$ for silicon--silicon pairs in the no-crossing-requirement
samples with QA cuts. Black: peak window $|m-\mu|<2\sigma$; red: sidebands 5--7$\sigma$ (same total
width); blue points: difference. Left \Ks, right \Lam. The subtracted signal sits at $X_1=X_2$.}
\label{fig:dX}
\end{figure}

\begin{figure}[H]
\centering
\includegraphics[width=0.95\textwidth]{fig_xclass_Ks}
\caption{$\pi^+\pi^-$ mass per crossing class A--F (no crossing requirement, QA cuts); blue: fit with
mean and width fixed to the default-sample peak. Numbers in Table~\ref{tab:xclass}.}
\label{fig:xclassKs}
\end{figure}
\begin{figure}[H]
\centering
\includegraphics[width=0.95\textwidth]{fig_xclass_Lambda}
\caption{As Figure~\ref{fig:xclassKs} for $p\pi$.}
\label{fig:xclassLambda}
\end{figure}

\section{\Dzero{} candidates}
\label{sec:d0}
The \Dzero{} ntuple holds \nCandDzero{} candidates (\candPerEvDzero{} per event) with the pre-cuts of
Tables~\ref{tab:trackcuts}--\ref{tab:precuts}; the stored variables are shown in
Figure~\ref{fig:precutD0}. Because only one $K\pi$ assignment per pair is kept
(Section~\ref{sec:reco}), Figure~\ref{fig:massD0} (left) also shows the mass recomputed with the
kaon and pion masses swapped. A tighter illustrative selection (DIRA$>0.95$, decay-length
significance $>3$, DCA$_{12}<100\,\mu$m, both daughters DCA$_{xy}>50\,\mu$m and with INTT clusters)
leaves \nDzeroSel{} candidates and no visible peak (Figure~\ref{fig:massD0}, right).

\caveat{No \Dzero{} signal is seen in \nEvents{} events. This is not unexpected for an unoptimised
selection on one run, and it does not test the \Dzero{} configuration. A meaningful check needs a
much larger sample and an optimised selection (e.g.\ using simulation).}

\begin{figure}[H]
\centering
\includegraphics[width=\textwidth]{fig_precut_D0}
\caption{\Dzero{} candidates as stored (\nCandDzero{} candidates): DIRA, DCA$_{12}$, SV $\chi^2$/ndf,
decay length, smaller daughter DCA$_{xy}$, \pt. Red dashed lines: pre-cuts (DIRA$>0$,
DCA$_{12}<0.05$\,cm, $\chi^2$/ndf$<20$).}
\label{fig:precutD0}
\end{figure}

\begin{figure}[H]
\centering
\includegraphics[width=0.9\textwidth]{fig_mass_D0}
\caption{$K\pi$ mass. Left: all stored candidates with the stored mass assignment (black) and with
$K\leftrightarrow\pi$ swapped (red). Right: illustrative tighter selection given in the panel
(\nDzeroSel{} candidates). The \Dzero{} mass is 1.865\,GeV/$c^2$; no peak is visible.}
\label{fig:massD0}
\end{figure}

\section{Running the analysis and proposed next steps}
\label{sec:production}

\paragraph{Commands} (in \fpath{/gpfs02/sphenix/user/dongxs/run3pp}):
\begin{itemize}\setlength\itemsep{0pt}
\item build the companion module once: \fcmd{./build.sh};
\item one segment interactively:\\ \fcmd{macro/run_kfp.sh <nEvents|-1> <DST or list> <outdir> [crossingStudy 0/1] [nSkip] [channels]};
\item a list on condor: \fcmd{condor/make_list.sh <run> <firstSeg> <nFiles> > list.txt}, then
  \fcmd{condor/submit.sh <tag> list.txt [crossingStudy] [loose|prod]} (outputs in \fpath{output/<tag>/});
\item merge: \fcmd{hadd -k} per channel; numbers and figures of this note:
  \fcmd{root -b -q 'analysis_note/scripts/make_numbers.C+("<merged dir>","<note dir>")'}.
\end{itemize}

\paragraph{Offline analysis checklist.} Remove duplicates with \fcmd{Collision_BCO}
(Section~\ref{sec:dup}); normalise with \fcmd{EventInfo}; choose the trigger with the
\fcmd{trigger_*} branches or \fcmd{GL1_scaled}; apply a PV-$z$ cut with \fcmd{primary_vertex_z};
compute the swapped \Dzero{} mass; use \fcmd{track_N_*_nHits} for daughter quality and
\fcmd{track_N_calculated_dEdx} for PID.

\decision{Crossing study channels: keep \fcmd{crossingStudy=0} for production (saves
\msPerEvKsNoX+\msPerEvLambdaNoX\,ms per event and $\approx$1.5\,GB per 30k events), and repeat the study
once on $\sim$300k events to settle the \Lam{} class-C question (Section~\ref{sec:crossing}).}

\choice{Ntuple size for production: resolved by the \fcmd{prod} cut set and slim format of
Section~\ref{sec:trim} (\trimKBDira\,kB per event, $\approx$\trimTBDira\,TB for the full dataset).}

\decision{\Dzero{} production cuts (daughter DCA$>30\,\mu$m, DIRA$>0.9$) were chosen without a visible
signal; their efficiency should be checked with simulation.}

\decision{Event selection: which trigger(s) and PV-$z$ range define the analysis sample (stored, not
applied).}

\section{Trimming the output for mass production}
\label{sec:trim}

The loose configuration of Section~\ref{sec:reco} stays the default for small diagnostic samples.
For the full-dataset production a second cut set, \fcmd{prod}, trims the output in four steps. It is
selected by the last argument of the macro (\fcmd{Fun4All_run3pp_KFP.C(...,"prod")}, i.e.\
\fcmd{macro/run_kfp.sh <nev> <dst> <outdir> 0 0 Ks,Lambda,D0 prod} or
\fcmd{condor/submit.sh <tag> <list> 0 prod}) and is always run without the crossing-study
channels. Each step was tested on the same 31 segments of run 79509 (\trimNev{} events).

\subsection{Trimming steps}
\begin{enumerate}\setlength\itemsep{2pt}
\item \textbf{No crossing-study output.} The \fcmd{*_noXmatch} files (\trimNoXMB\,MB, \trimNoXKB\,kB per
  event in the test) are not produced.
\item \textbf{Daughter impact parameter} (KFParticle \fcmd{setMinimumTrackPV_DCA}, 3D DCA, minimum over
  the PVs of the track's crossing, Section~\ref{sec:trackcuts}): $>0.05$\,cm for \Ks{} and \Lam{}
  daughters, $>0.003$\,cm (30\,$\mu$m) for \Dzero{} daughters; in this step also DIRA$>0.8$ for \Dzero.
\item \textbf{Slim format} (\fpath{tools/slim_kfp.C}, run inside the job after \fcmd{se->End()}): the
  64 trigger branches are not booked (\fcmd{getTriggerInfo(false)}, no \fcmd{TriggerRunInfoReco}; the
  trigger vectors remain per event in \fcmd{EventInfo}); redundant momentum branches are dropped
  (mother: $p_x,p_y,p_z,E,p,\sigma_p,\theta,y$, charge; daughters: $p_x,p_y,p_z,E,p,\sigma_p,\theta,y$,
  mass, $j_T$, $x,y,z$; and \fcmd{secondary_vertex_mass_pionPID}); the tree is rewritten with ZSTD
  compression and a 32\,kB basket size. The latter fixes a storage inefficiency of KFParticle, which
  books the covariance branches with a buffer of 6 or 21 \emph{bytes}, i.e.\ one basket per entry.
  For \Ks{} the four covariance branches took 640 bytes per candidate before and 260 after the fix
  (PV 121$\to$15, mother 171$\to$81, daughters 348$\to$163; \fpath{logs/branch_sizes_Ks.txt}).
\item \textbf{No covariance matrices} (mother 21, daughters $2\times21$, PV 6 elements).
\item \textbf{DIRA$>0.9$} for all three channels.
\end{enumerate}
The dropped kinematic quantities follow from the kept ones: $p=p_T\cosh\eta$, and the daughter mass
from its PDG code (needed e.g.\ for the swapped \Dzero{} mass). Duplicates across events are removed
with the key (\fcmd{Collision_BCO}, mass, \fcmd{track_1_pT}), because \fcmd{track_1_px} no longer exists.

\subsection{Size and signal}
\begin{table}[H]
\centering\footnotesize
\setlength\tabcolsep{4pt}
\caption{Cumulative effect of the trimming steps on the same \trimNev{} events
(\fpath{scripts/trim_study.C}). Candidates and sizes of the \Ks, \Lam, \Dzero{} trees; MB and kB per
event are file sizes of the merged ntuples (the \fcmd{*_noXmatch} files and the \fcmd{EventInfo}
file, \trimEvInfoKB\,kB/event in all versions, are not included). Signal: \Ks{} and \Lam+\aLam{}
yield with the offline QA cuts of Section~\ref{sec:qacuts} (Gaussian+pol2 fit, duplicates kept,
statistical errors).}
\label{tab:trim}
\begin{tabular}{rp{0.26\textwidth}rrrrrrrr}
\toprule
 & version & br. & \Ks{} cand. & \Lam{} cand. & \Dzero{} cand. & MB & kB/ev & \Ks{} signal & \Lam{} signal \\
\midrule
\input{tables/tab_trim.tex}
\bottomrule
\end{tabular}
\end{table}

\begin{itemize}\setlength\itemsep{1pt}
\item The three candidate trees shrink from \trimKBLoose{} to \trimKBDira\,kB per event (\trimFracDira\%
  of the loose size); with the crossing-study files of the diagnostic run included the reduction is a
  factor $\approx$\trimTotalReduction.
\item The largest single gains come from the daughter impact-parameter cut (version 0$\to$1:
  \trimFracProd\% of the loose size) and from the covariance branches (versions 1$\to$3: basket fix and
  removal, \trimKBProd$\to$\trimKBNoCov\,kB/event). The trigger block was only 0.6\% of the size.
\item With the QA cuts the \Ks{} and \Lam{} yields are unchanged within less than 1\% through all steps.
  Since the QA cuts already contain DIRA$>0.995$ and $|{\rm DCA}_{xy}|>0.02$\,cm, this shows that the
  trimming does not cut into that selection; it does not measure the loss for a looser final selection.
\item Signal removed by DIRA$>0.9$: with the other QA cuts but no DIRA cut, sideband-subtracted counts in
  the version before the change give \Ks{} $\bandLoKs\pm\ebandLoKs$ in $0.8<{\rm DIRA}\leq0.9$ against
  $\bandHiKs\pm\ebandHiKs$ above 0.9, and \Lam{} $\bandLoLambda\pm\ebandLoLambda$ against
  $\bandHiLambda\pm\ebandHiLambda$, i.e.\ no significant loss.
\end{itemize}

\caveat{The sizes in Table~\ref{tab:trim} are of merged files written by \fcmd{hadd} (zlib level 1).
The per-segment job outputs, which are what a production stores, measure 3.72, 1.76 and 0.62\,kB/event
for versions 1, 2 and 4 (\fpath{logs/trim_persegment.txt}), within 2--7\% of the merged values. The signal
loss of the impact-parameter cuts was not measured without the QA $|{\rm DCA}_{xy}|>0.02$\,cm cut, and
the \Dzero{} loss cannot be measured in this sample (no peak).}

\paragraph{Production estimate.} With \trimKBDira+\trimEvInfoKB\,kB per event, the $\approx5\times10^8$
events in the DST area at the time of the test (\nDSTfiles{} segments) give $\approx$\trimTBDira\,TB,
compared with \trimTBLoose\,TB for the loose configuration without the crossing-study files. The three
KFParticle modules need 63+71+94\,ms per event in the final configuration
(\fpath{logs/trim_persegment.txt}).

\section{Audit}
\label{sec:audit}

\subsection{Checks that passed}
\begin{itemize}\setlength\itemsep{1pt}
\item All jobs finished: \nCondorOK/\nCondorJobs{} condor jobs with exit code 0 and \nKFPdone/\nKFPlogs{}
  job logs ending with ``All done'' (\checkLogsDone; \fpath{scripts/job_stats.py} on
  \fpath{logs/condor_471173.txt}, \fpath{logs/kfp_job_summary.txt}).
\item KFParticle's own candidate count for \Ks, summed over the 31 job logs (\nCandKsFromLogs), equals the
  number of entries in the merged ntuple (\nCandKs) (\checkCandLogs; \fpath{scripts/job_stats.py}).
\item Crossing consistency of the default \Ks{} and \Lam{} samples: $X_1=X_2=X_{\rm PV}$ for every
  candidate (\checkXKs, \checkXLambda; \fpath{scripts/make_numbers.C}, full test sample).
\item Track crossing of TPC-only tracks is 0 for all 17871 TPC-only tracks of segment 79509-00000
  (\fpath{tools/inspect_dst.C}, \fpath{logs/inspect_dst_79509_00000.txt}).
\item Merged ntuples copied to the laptop are identical to the cluster files (md5).
\item Every figure was opened and compared with its caption.
\end{itemize}

\subsection{Not verified (assumptions)}
\begin{itemize}\setlength\itemsep{1pt}
\item The internal KFParticle behaviour described in Section~\ref{sec:reco} (order of cuts, PV choice,
  implicit ``vertex in own crossing'' track requirement, one candidate per pair) was read from the
  source of the coresoftware commit, not tested with dedicated inputs.
\item Interpretation of the vertex spike at $X=0$ as the triggered collision.
\item How the crossing of MVTX-only tracks is assigned in the production.
\item The trigger-bit names come from \fcmd{TriggerRunInfoReco}; their correctness for this run and the
  trigger mix of the sample were not checked.
\item The offline QA cuts are ad hoc; yields and S/B depend on them.
\item Fit model (Gaussian + pol2) and fit ranges: no systematic variation.
\item The duplicate key (rounded mass and $p_x$) could in principle merge two different candidates of
  the same collision; the effect was not quantified.
\item Data quality of run 79509.
\end{itemize}

\section{Open issues}
\label{sec:open}
Ordered by importance; [decision] marks items for the analyser.
\begin{enumerate}\setlength\itemsep{1pt}
\item \Lam{} with daughters in different crossings: $\xclCLambda\pm\exclCLambda$ in class C; repeat
  with $\sim$10$\times$ the statistics (Section~\ref{sec:crossing}).
\item \Ks{} mass $\massShiftKs$\,MeV below the PDG value: check momentum scale, field map and
  energy-loss treatment (Section~\ref{sec:v0}).
\item Production trimming (Section~\ref{sec:trim}): signal loss of the impact-parameter cuts for a final
  selection looser than the QA cuts, and \Dzero{} efficiency, still to be measured.
\item V0s with a TPC-only daughter are only recovered from the crossing-0 collision; recovering the
  others would need a crossing estimate for TPC-only tracks (not available in the DST).
\item [decision] \Dzero{} pre-cuts; no signal seen (Section~\ref{sec:d0}).
\item [decision] Event selection: trigger and PV-$z$ cut.
\item Include runs 80906 and 81600 (single-segment tests still running) and more runs to check
  run dependence.
\item Efficiency and acceptance (simulation with the same macro).
\end{enumerate}

\section{Glossary}
\label{sec:glossary}
\begin{description}\setlength\itemsep{0pt}
\item[BCO] beam-clock counter; \fcmd{GL1_BCO} is the BCO of the event's trigger.
\item[crossing ($X$)] bunch crossing of a track or vertex relative to the GL1 BCO.
\item[DCA$_{12}$] 3D distance of closest approach between the two daughter trajectories.
\item[DIRA] cosine of the angle between the mother momentum and the PV$\to$SV vector.
\item[FD $\chi^2$] flight-distance $\chi^2$ of the SV from the PV.
\item[I / M / T] daughter with INTT clusters / MVTX but no INTT / TPC-only.
\item[PV, SV] primary vertex (\fcmd{SvtxVertexMap}), secondary (decay) vertex.
\item[pre-cuts] requirements applied in KFParticle when the ntuple is written.
\item[QA cuts] offline cuts used in this note only to show the peaks (Section~\ref{sec:qacuts}).
\item[streaming readout] continuous TPC/silicon readout: one event contains collisions from many crossings.
\end{description}


\end{document}
